%%
%% Chapter 7: Discussion
%%
\chapter{Discussion}
\label{ch:07}

The preceding two chapters reported the fusion, backbone, robustness, and
test-time-adaptation results largely on their own terms. This chapter steps
back and discusses the patterns that cut across those sections: why a weak
skeleton stream still helps a strong visual-language stream, why the two
benchmarks prefer different visual backbones, why fixed-weight late fusion
fails where learned gating succeeds, what the failure of entropy-minimization
test-time adaptation reveals about LayerNorm-based fusion heads, and what the
frozen-backbone design costs and buys computationally. The chapter closes with
a full comparison against the published state of the art
(Section~\ref{sec:sota}), including the fair-subset analysis that defines the
regime in which this work's numbers are directly comparable.

\section{The Complementarity Argument}
\label{sec:disc-complementarity}

Despite achieving only 68.8\% AUC in isolation on UCF-Crime (and 40.8\% AP on
XD-Violence), the skeleton modality consistently improves gated fusion
performance over visual-only baselines. Across all eight
backbone$\times$dataset configurations evaluated in Chapter~\ref{ch:05}, gated
fusion meets or exceeds its visual-only counterpart --- strictly positive in
six configurations with two ties, a pattern unlikely under a no-effect null
(tie-dropping sign test, $p \approx 0.03$) --- with a mean gain of $+0.6$
points. The gain is small but consistent, and that combination is precisely
what one expects from a modality that contributes \emph{orthogonal}
information: motion dynamics that appearance-based vision-language models do
not capture, rather than a redundant re-encoding of what the visual stream
already knows.

The strongest single piece of evidence that the network extracts genuine
signal from the skeleton stream is the gated mechanism's large advantage over
fixed equal-weight late fusion on XD-Violence, where gating raises the CLIP
backbone's AP from 63.9\% to 76.5\% --- a 12.6-point gap. A redundant or
pure-noise modality could not produce this behavior: if the skeleton features
carried no usable information, no input-dependent weighting scheme could
recover 12 points of AP by consulting them. The gap instead indicates that the
skeleton stream is informative for some inputs and misleading for others, and
that the value of the modality is only realized when the network can decide,
per input, how much to trust it.

The per-class analysis in Appendix~\ref{app:b} probes the mechanism directly.
If skeleton features contribute motion information, the fusion gain should
concentrate in human-motion-centric anomaly categories. The observed pattern
is directionally supportive but mixed and noisy: human-motion categories show
a mean per-class gain of $+0.31 \pm 1.08$ points (positive in 7 of 10
categories), appearance-centric categories $-0.29 \pm 0.16$ (positive in 0 of
3), a difference of $+0.60$ points, with Shooting ($+1.05 \pm 0.20$) the
cleanest single-category signal. Per-class sample sizes on UCF-Crime are
small, so we present this as a consistency check on the mechanism rather than
as an independent statistical result.

Two honest qualifications frame the argument. First, the magnitude is modest
and backbone-dependent: complementarity here means ``meets or exceeds, usually
by a little,'' not a transformative gain. Second, the complementarity claim is
about the \emph{fusion} of the streams, not about parity between them --- the
skeleton stream on its own is far below the visual stream on both benchmarks,
and nothing in our results suggests otherwise.

\section{Dataset-Specific Backbone Preferences}
\label{sec:disc-backbones}

The divergence between UCF-Crime and XD-Violence in optimal backbone selection
--- SigLIP2 Giant for UCF-Crime ($82.5 \pm 0.4$\% AUC under gated fusion),
SigLIP2 SO400M for XD-Violence ($78.7 \pm 0.9$\% AP) --- has practical
implications for system deployment. UCF-Crime consists of fixed-camera
surveillance footage where the larger Giant backbone's richer representations
may better disambiguate fine-grained anomaly categories. XD-Violence draws
from heterogeneous sources including movies, sports broadcasts, and news
footage, where SO400M's shape-optimized architecture may provide a better
capacity fit.

Notably, the XD-Violence preference is not visible in the visual-only
baselines, which sit within seed variance of one another; it emerges only
under fusion. This means the interaction is not simply ``a better visual
encoder wins'' --- the preference reflects how each backbone's feature
geometry combines with the skeleton stream through the learned gate. A
backbone ranking obtained from visual-only probes, or from general-purpose
retrieval and classification benchmarks, would not have predicted the fusion
ranking on XD-Violence.

This finding cautions against selecting a visual backbone for anomaly
detection based on general-purpose benchmarks alone and suggests that
dataset-specific evaluation is essential. The frozen-backbone design makes
such evaluation cheap (Section~\ref{sec:disc-efficiency}): swapping the
backbone requires re-extracting features once and retraining only the fusion
head, so a per-deployment backbone sweep is a matter of minutes of training
time per candidate rather than days of fine-tuning.

\section{The Late-Fusion Failure Mode}
\label{sec:disc-latefusion}

The consistent underperformance of late fusion relative to both gated fusion
and, in some cases, visual-only baselines warrants attention. Late fusion
combines the two modalities with a fixed equal weight, so when the skeleton
signal quality is low --- as reflected by its 68.8\% standalone AUC on
UCF-Crime and 40.8\% AP on XD-Violence --- it still contributes half of every
score, degrading inputs where the visual stream alone is more reliable. On
UCF-Crime the best late-fusion configuration reaches only $79.9 \pm 0.8$\%
AUC, and the CLIP late-fusion variant (78.9\%) falls below its own visual-only
baseline (81.2\%): adding the skeleton stream at a fixed mix actively hurts.
On XD-Violence the failure is more dramatic --- late fusion collapses to
63.6--65.7\% AP across backbones, more than 10 points below gated fusion in
every configuration.

The asymmetry between the two datasets is informative. The quality gap between
the modalities is much larger on XD-Violence (skeleton-only 40.8\% AP versus
visual-only in the mid-70s) than on UCF-Crime, and equal-weight averaging is
most damaging precisely when the mixed streams are most unequal. The gated
mechanism instead learns an input-dependent, per-dimension reweighting over a
residual baseline, allowing it to reduce the skeleton stream's relative
contribution for inputs where visual features are more informative --- and to
do so continuously, per feature dimension, rather than with a single global
mixing coefficient. The lesson for multi-modal VAD systems is that the
decision of \emph{how much} to trust an auxiliary modality is itself learnable
signal, and discarding it (as any fixed mixing scheme does) can cost more than
the auxiliary modality contributes.

\section{Test-Time Adaptation and the LayerNorm Barrier}
\label{sec:disc-ln-barrier}

Our TTA results (Chapter~\ref{ch:06}) are an informative negative finding with
implications beyond VAD. Entropy-minimization methods such as
TENT~\cite{wang2021tent} and SAR~\cite{niu2023sar} were designed for BatchNorm
networks, where adaptation recalibrates the running mean and variance
statistics that summarize the source distribution~\cite{ioffe2015batchnorm}.
Our fusion head normalizes with LayerNorm~\cite{ba2016layernorm}, which keeps
no such running statistics --- normalization is computed per sample --- so on
each batch these methods can only nudge the LayerNorm affine parameters
($\gamma$, $\beta$; 1{,}536 scalars in our head). With the visual-language
backbone, the gate, and the detection head all frozen, the corruption-induced
shift lives upstream in the frozen features, and affine updates of this size
are far too small to reorder a rank-based AUC --- leaving the metric
essentially unchanged ($<0.1$ points across all four backbones and three
seeds).

The failure is structural, not a matter of tuning. A rank-based metric such as
AUC is invariant to any monotone transformation of the scores; to improve it,
an adaptation method must \emph{reorder} snippets, and a 1,536-parameter
affine correction applied after frozen feature extraction has almost no
capacity to do so. Forcing the features back toward the clean statistics
directly, in the style of statistic-restoration methods such as
NORM~\cite{schneider2020norm} or correlation alignment~\cite{sun2016coral},
does reach the right space but is rank-disruptive in our experiments, with an
effect whose sign flips across backbones. The lesson is not that the
corruptions are benign --- source-only corrupted AUC drops by roughly twenty
points --- but that adaptation confined to LayerNorm affine parameters cannot
move a frozen, rank-based detector.

What does work here is to change \emph{which} modality the frozen fusion
trusts: the label-free discriminative-reliability reweighting of
Section~\ref{sec:disc} routes the gated fusion toward the corruption-robust
skeleton stream when the visual-language stream's discriminative signal
collapses, recovering corrupted-AUC on all four backbones (mean $+1.21$
points) without updating a single parameter. This suggests that robustness for
frozen LayerNorm-based fusion heads is better pursued through
reliability-aware routing --- or, more broadly, prompt- or adapter-based
mechanisms that alter the features themselves --- than through affine
test-time adaptation. We emphasize that these gains concern corrupted-input
robustness only, are largest where exactly one modality degrades, and rely on
a transductive protocol whose limitations are discussed in
Chapter~\ref{ch:08}.

\section{Computational Efficiency}
\label{sec:disc-efficiency}

Our frozen-backbone design yields a lightweight trainable component: the gated
fusion head contains 0.50M (CLIP) to 1.02M (SigLIP2 Giant) trainable
parameters, scaling with the visual feature dimension $d_v$, as only the
projection, gating, and MIL detection-head layers are learned. Training
completes in under 5 minutes per configuration on a single RTX 4090. This
efficiency enables the comprehensive multi-backbone evaluation presented in
this thesis --- four visual backbones, two datasets, six fusion variants, and
three seeds per configuration --- which would be infeasible with end-to-end
backbone fine-tuning on the same hardware.

The efficiency argument is a deployment argument as much as a research-budget
one. Because all learned capacity lives in a sub-megaparameter head on top of
cached features, the marginal cost of retraining --- for a new site, a new
threshold policy, or a newly released vision-language backbone --- is minutes,
not GPU-days, and the feature extraction pipeline is shared across every
variant. Section~\ref{sec:sota} returns to this point quantitatively when
comparing against methods whose inference alone requires data-center-class
GPUs.

\section{Comparison with the State of the Art}
\label{sec:sota}

\noindent\textit{Metric definitions (stated once).}
On \textbf{UCF-Crime} we report \textbf{frame-level ROC-AUC (\%)}; on
\textbf{XD-Violence} we report \textbf{Average Precision (AP) / AUPRC (\%)}
--- \emph{not} ROC-AUC. Several cited papers publish both an XD ROC-AUC and an
XD AP (e.g.\ EventVAD~\cite{shao2025eventvad} 87.51 AUC vs 64.04 AP;
LAVAD~\cite{zanella2024lavad} 85.36 AUC vs 62.01 AP); we
tabulate only the AP. The \textbf{Setting} column encodes, compactly, whether
the backbone is frozen vs.\ (partly) fine-tuned, whether a learnable text/prompt
branch is used, whether audio is used, and whether the reported metric is
frame-level or video-level, so a reader can see at a glance which rows are
\emph{not} apples-to-apples with a frozen, visual-only, frame-level system.
We do \emph{not} claim state-of-the-art performance; the honest claim is
competitiveness within the constrained frozen / visual-only / single-GPU regime
(Table~\ref{tab:sota-fair}).

\begin{table}[p]
\centering
\singlespacing
\caption{Published-numbers comparison on UCF-Crime (frame-level ROC-AUC, \%)
and XD-Violence (Average Precision, \%), sorted by UCF AUC descending. The
\textbf{This work} row is placed at its true position and shown in
\textbf{bold}; it is never a bare head-to-head, because the \textbf{Setting}
column makes each row's regime explicit. Superscripts mark per-number caveats
(see notes below the table).}
\label{tab:sota-full}
\small
\setlength{\tabcolsep}{4pt}
\resizebox{\textwidth}{!}{%
\begin{tabular}{l c l c c l}
\toprule
Method & Year & Venue & UCF AUC & XD AP & Setting \\
\midrule
GS-MoE$^{a}$~\cite{damicantonio2025gsmoe} & 2025 & ICCV'25 (arXiv)   & 91.58 & 82.89 & Frozen feats, no-text, frame; MoE experts \\ % VERIFIED 2026-06-22: GS-MoE 91.58 UCF / 82.89 XD AP confirmed from arXiv:2508.06318 Table 1 (I3D, RGB-only, no text); ICCV 2025 poster
PI-VAD~\cite{majhi2025pivad}              & 2025 & CVPR'25           & 90.33 & 85.37 & Frozen I3D + 5 train-only modalities, text(train), frame \\
Holmes-VAD$^{b}$~\cite{zhang2024holmesvad} & 2024 & arXiv'24          & 89.51 & 90.67 & MLLM LoRA-tuned, text, frame; heavy VRAM \\ % RE-VERIFY: Holmes-VAD 89.51/90.67 arXiv preprint, no peer-reviewed venue; partial MLLM LoRA training (not frozen-feature)
DSANet~\cite{yin2026dsanet}               & 2026 & AAAI'26           & 89.44 & 86.95 & Frozen CLIP + V-L alignment, text, frame \\
Holmes-VAU$^{c}$~\cite{zhang2025holmesvau} & 2025 & CVPR'25           & 88.96 & 87.68 & LoRA MLLM + sampler, text, \emph{VIDEO-level} \\ % RE-VERIFY: Holmes-VAU 88.96/87.68 are VIDEO-LEVEL (398-sample split), NOT frame-level; disclose or omit from frame-level head-to-head
STPrompt~\cite{wu2024stprompt}            & 2024 & ACM MM'24         & 88.08 & N/A$^{d}$ & Frozen CLIP + ST prompts, text, frame; XD not eval. \\
FDPN$^{e}$~\cite{song2025fdpn}            & 2025 & WACV'25           & 88.03 & N/A$^{e}$ & I3D RGB, frame; self-report, XD not eval. \\ % RE-VERIFY: FDPN 88.03 LOW conf, niche WACV-25 paper's OWN method (+0.01 over VadCLIP self-report); XD not evaluated -- re-check primary or drop
VadCLIP~\cite{wu2024vadclip}              & 2024 & AAAI'24           & 88.02 & 84.51 & Frozen CLIP + dual-branch, text, frame \\
TPWNG~\cite{yang2024tpwng}                & 2024 & CVPR'24           & 87.79 & 83.68 & CLIP text-proj fine-tuned, text, frame; self-train \\
CLIP-TSA~\cite{joo2023cliptsa}            & 2023 & ICIP'23           & 87.58 & 82.19$^{f}$ & Frozen CLIP \emph{visual-only}, no-text, frame \\ % RE-VERIFY: CLIP-TSA XD = 82.19 (primary AUC@PR); VadCLIP table = 82.17 (0.02 drift); NEVER the paper's own 94.02 (non-comparable protocol)
MGFN~\cite{chen2023mgfn}                  & 2023 & AAAI'23           & 86.98 & 79.19$^{g}$ & Frozen I3D feats, no-text, frame \\ % RE-VERIFY: MGFN XD = 79.19 is the I3D-RGB AP (NOT 80.11, which is the VideoSwin variant); keep one backbone (I3D)
UR-DMU~\cite{zhou2023urdmu}               & 2023 & AAAI'23           & 86.97 & 81.66 & Frozen I3D feats, no-text, frame; single-GPU \\
PEL4VAD~\cite{pu2024pel4vad}              & 2024 & IEEE TIP'24       & 86.76 & 85.59 & Frozen I3D + CLIP text prompts, text, frame \\
PiercingEye$^{h}$~\cite{leng2026piercingeye} & 2026 & IEEE TPAMI'26     & 86.64 & 88.82 & Frozen I3D+VGGish+text, \emph{audio}+text, frame \\ % RE-VERIFY: PiercingEye XD = 88.82 LOW conf (inferred from relative gains) and is AUDIO+visual+text; UCF 86.64 is visual-only
AnomalyCLIP$^{i}$~\cite{zanella2024anomalyclip} & 2024 & CVIU'24           & 86.36 & 78.51 & Frozen CLIP latent dirs, text, frame \\ % RE-VERIFY: AnomalyCLIP 86.36/78.51 LOW conf (not re-verified from primary table); the ~90.3 leaderboard figure is recognition mAUC, a DIFFERENT task -- do not cite it as detection AUC
TEVAD$^{j}$~\cite{chen2023tevad}          & 2023 & CVPRW'23          & 84.9  & 79.8  & Frozen I3D + caption text, text, frame \\ % RE-VERIFY: TEVAD XD ~= 79.8 MEDIUM precision (headline 79.8, ablation variants 79.3-79.76); disregard '88.28' web misattribution
Light-WVAD~\cite{wang2024lightwvad}       & 2024 & Neurocomp.'24     & 84.7  & ---   & Frozen I3D, no-text, frame; consumer-GPU head (XD not evaluated) \\
RTFM~\cite{tian2021rtfm}                  & 2021 & ICCV'21           & 84.30 & 77.81 & Frozen I3D feats, no-text, frame; lightweight \\
EventVAD$^{k}$~\cite{shao2025eventvad}    & 2025 & ACM MM'25         & 82.03 & 64.04 & Training-free 7B MLLM, text, frame; 80GB A800 \\ % RE-VERIFY: EventVAD XD = 64.04 (AP); do NOT cite 87.51 (that is the XD ROC-AUC). Training-free != low-compute
\textbf{This work (dual-modal fusion)} & \textbf{2026} & \textbf{(thesis)} & \textbf{82.5} & \textbf{78.7} & \textbf{Frozen skel+CLIP/SigLIP2, no-text, frame; 1$\times$RTX 4090} \\
LAVAD~\cite{zanella2024lavad}             & 2024 & CVPR'24           & 80.28 & 62.01 & Training-free 13B-LLM pipeline, text, frame \\ % RE-VERIFY: LAVAD XD = 62.01 (AP), NOT 85.36 (its XD ROC-AUC); double-check AP not AUC
Sultani et al.$^{l}$~\cite{sultani2018ucfcrime} & 2018 & CVPR'18           & 75.41 & ---$^{l}$ & Frozen C3D, FC+MIL, no-text, frame; founding baseline \\ % RE-VERIFY: Sultani XD 73.20 is NOT from the 2018 paper -- it is Wu et al. ECCV 2020's re-impl; attribute correctly or omit the XD cell
HyperVD$^{m}$~\cite{peng2024hypervd}      & 2024 & Image\,\&\,Vis.\,Comp.'24 & N/A$^{m}$ & 85.67 & Frozen I3D+VGGish, \emph{audio}, frame; UCF not rep. \\ % RE-VERIFY: HyperVD is AUDIO-VISUAL (85.67 AP); visual-only variant = 82.51; UCF not reported -- label modality if compared
Ghadiya et al.$^{n}$~\cite{ghadiya2024crossmodal} & 2024 & CVPRW'24 (MULA)   & N/A$^{n}$ & 86.34 & Frozen I3D+VGGish, \emph{audio}, frame; UCF not rep. \\ % RE-VERIFY: Ghadiya 86.34 MEDIUM conf, AUDIO-VISUAL workshop preprint; UCF not reported -- not comparable to visual/skeleton-only
\bottomrule
\end{tabular}%
}

\vspace{2pt}
\scriptsize
\begin{flushleft}
\textit{Metric:} UCF AUC $=$ frame-level ROC-AUC (\%); XD AP $=$ Average
Precision / AUPRC (\%). \textbf{We do not claim state-of-the-art}; rows above
this work add a learnable text branch, audio, an MLLM, partial fine-tuning, or
a video-level metric (see Setting / notes), so a bare head-to-head against them
is not apples-to-apples (Table~\ref{tab:sota-fair}).\par\smallskip
\textit{Per-number caveats.}
$^{a}$\,GS-MoE 91.58\,/\,82.89: confirmed from the arXiv:2508.06318 primary table
(I3D, RGB-only, no text); ICCV 2025 poster. Omitted from the main comparison only
because its trained per-category mixture-of-experts head exceeds the single-GPU-class regime.
$^{b}$\,Holmes-VAD 89.51/90.67: arXiv preprint (no confirmed peer-reviewed venue);
MLLM LoRA fine-tuned --- not a frozen-feature head.
$^{c}$\,Holmes-VAU 88.96/87.68: \textbf{VIDEO-LEVEL} AUC/AP on a 398-sample split,
\emph{not} frame-level; apples-to-oranges vs.\ all frame-level rows including this
work --- do not tabulate without disclosing the metric difference.
$^{d}$\,STPrompt: XD-Violence is \emph{not} evaluated (benchmarks: UCF /
ShanghaiTech / UBnormal); any XD value would be fabricated.
$^{e}$\,FDPN 88.03: LOW confidence; the paper's own method, $+0.01$ over VadCLIP
(a ``narrowly beats'' self-report); XD not evaluated.
$^{f}$\,CLIP-TSA XD $=$ 82.19 (primary AUC@PR); the 82.17 that circulates is from
VadCLIP's secondary table; the paper's own 94.02 uses a non-comparable protocol
--- do not use it.
$^{g}$\,MGFN XD $=$ 79.19 is the \textbf{I3D-RGB} AP; 80.11 is the VideoSwin
variant. UCF 86.98 is I3D; we keep a single (I3D) backbone per row.
$^{h}$\,PiercingEye XD $=$ 88.82 is \textbf{audio$+$visual$+$text}; UCF 86.64 is
visual-only; LOW confidence (AP inferred/cross-checked).
$^{i}$\,AnomalyCLIP 86.36/78.51: LOW confidence (detection-task numbers); the
$\sim$90.3 leaderboard figure is recognition mAUC, a different task.
$^{j}$\,TEVAD XD $\approx$ 79.8: MEDIUM precision (ablation variants 79.3--79.76);
the ``88.28'' web snippet is a misattribution.
$^{k}$\,EventVAD XD $=$ 64.04 is the AP; the paper's XD ROC-AUC is 87.51 (different
metric); training-free but runs a 7B MLLM on an 80GB A800.
$^{l}$\,Sultani XD: the 73.20 AP often quoted is \emph{not} from the 2018 paper
(XD did not exist until 2020); it is Wu et al.'s ECCV-2020 re-implementation ---
attribute to Wu et al., not Sultani; XD cell omitted here.
$^{m}$\,HyperVD UCF not reported (XD-focused); XD 85.67 is audio-visual
(visual-only variant $=$ 82.51).
$^{n}$\,Ghadiya et al.\ 86.34: MEDIUM confidence, audio-visual workshop preprint;
UCF not reported --- not comparable to visual/skeleton-only methods.
\end{flushleft}
\end{table}

%% --- Fair-subset (apples-to-apples) table ---
\begin{table}[p]
\centering
\singlespacing
\caption{Fair-subset (apples-to-apples) comparison. \textbf{Subset definition:}
frozen features, no backbone fine-tuning, no learnable text/prompt branch,
visual-only RGB (or RGB$+$skeleton), and a single-GPU-class feature head. This
is the regime in which this work's 82.5\% UCF AUC / 78.7\% XD AP is most
directly comparable; methods with text/prompt branches, audio, MLLMs, partial
fine-tuning, or video-level evaluation are excluded (reasons in
Table~\ref{tab:sota-full}). Because \emph{every} row is, by construction, the
same setting, a per-row Setting column is unnecessary here --- the setting is
captured once by this caption. \textbf{Metric:} UCF AUC $=$ frame-level ROC-AUC
(\%); XD AP $=$ Average Precision (\%). \textbf{This work} in \textbf{bold}.}
\label{tab:sota-fair}
\small
\setlength{\tabcolsep}{5pt}
\begin{tabular}{l c c c p{0.40\linewidth}}
\toprule
Method & Year & UCF AUC & XD AP & Why it qualifies \\
\midrule
CLIP-TSA~\cite{joo2023cliptsa} & 2023 & 87.58 & 82.19 & Frozen CLIP \emph{visual-only}, no text branch --- the closest analogue (CLIP feats, no prompts) \\
UR-DMU~\cite{zhou2023urdmu}    & 2023 & 86.97 & 81.66 & I3D frozen feats, no text, no fine-tune, RGB-only \\
MGFN-I3D~\cite{chen2023mgfn}   & 2023 & 86.98 & 79.19 & I3D frozen feats, no text, RGB-only (I3D XD to keep one backbone) \\
Light-WVAD~\cite{wang2024lightwvad} & 2024 & 84.7  & ---   & Frozen I3D, lightweight head, no text --- explicit consumer-GPU efficiency play (XD not evaluated) \\
RTFM~\cite{tian2021rtfm}       & 2021 & 84.30 & 77.81 & Frozen I3D feats, no text, the canonical MIL anchor \\
\textbf{This work} & \textbf{2026} & \textbf{82.5} & \textbf{78.7} & Frozen skeleton (CTR-GCN) $+$ frozen CLIP/SigLIP2, gated fusion $+$ MIL head only, no text, single 4090 \\
Sultani et al.~\cite{sultani2018ucfcrime} & 2018 & 75.41 & ---   & Frozen C3D, FC $+$ MIL ranking; founding baseline (XD number is a later re-impl, omitted) \\
\bottomrule
\end{tabular}

\vspace{2pt}
\footnotesize
\begin{flushleft}
\textit{Read of the fair subset.} Within this strictly-frozen / no-text /
RGB($+$skeleton)-only band, this work's \textbf{78.7\% XD AP is essentially on
par with RTFM (77.81) and MGFN-I3D (79.19)} and clearly
above the Sultani lineage, while sitting below the CLIP-visual-feature ceiling
of CLIP-TSA (82.19) and UR-DMU (81.66). On UCF, 82.5\% trails the frozen-feature
pack (RTFM 84.30, UR-DMU 86.97, CLIP-TSA 87.58) but exceeds Sultani's 75.41
baseline by $\sim$7 points. The closest single comparator, CLIP-TSA, uses CLIP
\emph{visual} features with a temporal self-attention module and
multi-crop-style aggregation that this work deliberately omits; we hypothesize
the residual gap to it stems from that temporal-modeling and aggregation
machinery rather than from feature quality, but we do not isolate this
experimentally.
\end{flushleft}
\end{table}

\subsection{Positioning}
\label{sec:sota-positioning}

\noindent\textbf{Positioning.} The central claim of this work is a contribution
\emph{orthogonal} to peak benchmark AUC---a reproducible, single-GPU dual-modal
fusion pipeline with a label-free robustness mechanism; within-regime
competitiveness (clearest on XD-Violence) is presented below as supporting
evidence, not as a claim to the leaderboard.\par\medskip

\noindent\textbf{Scope and the deliberate-constraint frame.}
Our system reports 82.5\% frame-level AUC on UCF-Crime and 78.7\% AP on
XD-Violence under a deliberately constrained regime: both the skeleton backbone
(CTR-GCN, pre-trained on NTU-120) and the vision-language backbone
(CLIP/SigLIP2) are kept fully frozen, only a gated fusion module and an MIL
classifier are trained, and no learnable text-prompt branch, no audio stream,
and no heavy temporal graph are used, all on a single 24\,GB RTX 4090. These are
not incidental shortcomings but the experimental envelope of the thesis. The
strongest reported numbers in the field are obtained under materially heavier
budgets: VadCLIP~\cite{wu2024vadclip} (88.02\% UCF / 84.51\% XD, AAAI 2024),
TPWNG~\cite{yang2024tpwng} (87.79 / 83.68,
CVPR 2024) and DSANet~\cite{yin2026dsanet} (89.44 / 86.95, AAAI 2026) all add a
learnable vision-language text-alignment branch; PI-VAD~\cite{majhi2025pivad}
(90.33 / 85.37, CVPR 2025) induces
five auxiliary modalities at training time; and
Holmes-VAD~\cite{zhang2024holmesvad}/Holmes-VAU~\cite{zhang2025holmesvau}
fine-tune a multimodal LLM. The 5--9 point UCF gap to these systems is
therefore best read as the cost of the components we intentionally excluded,
not as a deficiency of the fusion approach itself.

\medskip\noindent\textbf{The contribution axis is orthogonal to peak AUC.}
The thesis does not claim a new state of the art on clean benchmark AUC; its
contributions lie on a different axis. First, it demonstrates a working
dual-modal fusion of frozen skeleton dynamics (CTR-GCN) with frozen
vision-language semantics through learned gated fusion, a combination not
represented among the visual-only and audio-visual methods above, and quantifies
the fusion gain over each single-stream baseline (skeleton-only 68.8\% / 40.8\%;
visual-only CLIP 81.2\% / 74.6\%; SigLIP2-Giant 82.4\% / 76.8\%). Second, it
provides a systematic four-backbone comparison (CLIP ViT-B/16, SigLIP2 ViT-B/16,
SigLIP2 SO400M, SigLIP2 Giant) revealing dataset-specific backbone preferences.
Third, and most distinctively, it studies test-time adaptation for frozen
dual-modal VAD: it shows that entropy-minimization TTA
(TENT~\cite{wang2021tent}, SAR~\cite{niu2023sar}) cannot
improve rank-based detection metrics, and proposes a label-free
discriminative-reliability reweighting that improves \emph{corrupted-AUC}
(robustness, not clean-benchmark AUC) on all four backbones (mean $+1.2$ points, up
to $+13.2$ points on the hardest corruption). None of the higher-AUC methods above
report a corruption-robustness or test-time-adaptation study of this kind, so
the work occupies a complementary niche rather than competing for the top
leaderboard slot.

\medskip\noindent\textbf{Competitiveness within the comparable regime (supporting evidence).}
With the contribution framed above, within-regime competitiveness is supporting
evidence rather than the central claim, and it differs by metric. Restricted to
the apples-to-apples subset (Table~\ref{tab:sota-fair}) --- fully frozen
features, no text branch, no fine-tuning, a single-GPU-class head --- on
XD-Violence our 78.7\% AP is on par with RTFM~\cite{tian2021rtfm} (77.81\%,
ICCV 2021), MGFN's~\cite{chen2023mgfn} I3D
variant (79.19\%), and well above
the founding Sultani et al.~\cite{sultani2018ucfcrime} MIL baseline. On
UCF-Crime, our 82.5\% clears
Sultani (75.41\%) but trails every modern method in this same frozen-feature
band --- RTFM (84.30\%), Light-WVAD~\cite{wang2024lightwvad} (84.7\%),
UR-DMU~\cite{zhou2023urdmu} (86.97\%), MGFN (86.98\%),
and CLIP-TSA~\cite{joo2023cliptsa} (87.58\%) --- so the competitiveness is
clearest on XD-Violence AP
and weakest on UCF-Crime AUC. The closest single comparator is CLIP-TSA
(87.58\% / 82.19\%), which, like us, uses frozen CLIP visual features with no
text branch; we hypothesize that the residual gap to it stems from its temporal
self-attention module and multi-crop aggregation --- components our minimal head
omits --- rather than from feature quality, but we do not isolate this
experimentally. Read against its true peer group rather than against fine-tuned
or multimodal systems, the fusion pipeline is a reasonable mid-pack entry on
XD-Violence rather than an outlier.

\medskip\noindent\textbf{Reproducibility and transparency as a deliberate virtue.}
All thesis numbers are reported as means over three random seeds on fixed
train/val splits with a single evaluation protocol, and the headline figures are
git-reproducible from a committed frame-count manifest. This contrasts with
several of the highest numbers in the literature, which are single-run preprints
(e.g.\ Holmes-VAD's~\cite{zhang2024holmesvad} 89.51\% / 90.67\%, an unrefereed
preprint).
It also contrasts with comparisons that quietly mix evaluation protocols ---
Holmes-VAU's~\cite{zhang2025holmesvau} 88.96\% / 87.68\% are video-level
metrics on a 398-sample split,
not the frame-level ROC-AUC / AP that every weakly-supervised baseline (and this
thesis) uses. By keeping the backbone frozen, the head tiny, the seeds fixed,
and the metric definitions explicit (UCF $=$ frame-level ROC-AUC; XD $=$ Average
Precision), the work prioritizes a number a reader can trust and reproduce over
a number that is merely higher.

\medskip\noindent\textbf{Compute honesty: training-free is not low-compute.}
A final positioning point concerns budget. Two recent training-free LLM methods,
LAVAD~\cite{zanella2024lavad} (80.28\% / 62.01\%, CVPR 2024) and
EventVAD~\cite{shao2025eventvad} (82.03\% / 64.04\%, ACM MM
2025), require running 13B and 7B multimodal LLMs at inference (LAVAD on dual
RTX 3090s, EventVAD on an 80\,GB A800), yet our frozen dual-modal head on a single
24\,GB consumer GPU matches or exceeds them on UCF and substantially exceeds both
on XD-Violence AP (78.7\% vs 62.01\% / 64.04\%). This makes a concrete point that
``training-free'' does not imply ``cheap,'' and that a small trainable head on
frozen features remains a strong, deployable trade-off for a single-GPU
researcher. The thesis's value is precisely this corner of the design space:
reproducible, consumer-GPU-feasible, multimodal-by-fusion, and explicitly
robustness-aware, rather than a pursuit of the absolute highest benchmark AUC.
